\section{Riemann--Siegel structure and Newman-time evolution}
\label{sec:RS}

The preceding acquisition theorems accept complete approximants with
paid errors. Riemann--Siegel theory supplies such approximants by
preserving oscillatory arithmetic terms and their remainder. It is
useful to separate three questions: where the classical saddles lie,
how the complete expression evolves in Newman time, and how its error
passes through an integrated logarithm. The first two identities alone
do not answer the signed upper-value question.

\subsection{The critical strip comes before the critical line}
For the standard oriented line \(\mathcal C\) from northeast to
southwest crossing the real interval \((0,1)\), with the usual
continuous branch of \(\log w\), Riemann's auxiliary function is
\begin{equation}\label{eq:RS-auxiliary}
 \mathscr R(s)=\int_{\mathcal C}
 \frac{w^{-s}\ee^{\pi\ii w^2}}{\ee^{\pi\ii w}-\ee^{-\pi\ii w}}\dd w.
\end{equation}
Its reflected identity is
\begin{equation}\label{eq:RS-reflection}
 \zeta(s)=\mathscr R(s)+\chi(s)\overline{\mathscr R(1-\bar s)},
 \qquad
 \chi(s)=\pi^{s-1/2}\frac{\Gamma((1-s)/2)}{\Gamma(s/2)}.
\end{equation}
These are the classical entire auxiliary integral and reflected zeta
formula; at exceptional factors they are understood by continuation
\cite[Section 3, Definition 4 and Theorem 5]{Arias2024AuxiliaryBasic}.
Contour displacement across integer poles produces finite Dirichlet
sums plus a remaining saddle contour. The corresponding two-sum
approximate functional equation is valid off the line as well. A
balanced cutoff near \(\sqrt{t/(2\pi)}\) does not require
\(\Re s=1/2\) \cite[Sections 2 and 8]{Arias2024AuxiliaryExpansions}.

On \(s=1/2+\ii t\), the two terms are related by conjugacy and
\(|\chi(s)|=1\). Choose the continuous gamma phase \(\theta(t)\) so
that \(\chi(1/2+\ii t)=\ee^{-2\ii\theta(t)}\). Multiplying by
\(\ee^{\ii\theta(t)}\) makes the full expression real, and its balanced
finite sums combine into
\begin{equation}\label{eq:RS-cosine}
 Z(t)=2\sum_{n\le\lfloor\sqrt{t/(2\pi)}\rfloor}
 n^{-1/2}\cos(\theta(t)-t\log n)+\mathcal E_{\rm RS}(t).
\end{equation}
The correction \(\mathcal E_{\rm RS}\) is retained. This is a
large-height expansion specialized by symmetry, not an expansion
assuming all zeros lie on the line. Its reality follows from the
functional equation whether or not RH holds. Away from the line the
two complete contributions must be kept separately.

\subsection{Off-line zeros do not force degenerate classical saddles}
To specify what is meant by a saddle, place the denominator of
\eqref{eq:RS-auxiliary} in the amplitude and choose the phases
\[
 \varphi_+(w;s)=\pi\ii w^2-s\log w,\qquad
 \varphi_-(w;s)=-\pi\ii w^2+(s-1)\log w.
\]
The second is the reflected phase in \eqref{eq:RS-reflection}.
For \(s=\sigma+\ii t\), \(t>0\), their near-positive saddle branches obey
\begin{equation}\label{eq:RS-saddles}
 w_+^2=\frac{t}{2\pi}-\frac{\ii\sigma}{2\pi},\qquad
 w_-^2=\frac{t}{2\pi}+\frac{\ii(1-\sigma)}{2\pi},\qquad
 \varphi_+''(w_+)=4\pi\ii,\quad \varphi_-''(w_-)=-4\pi\ii.
\end{equation}
These formulas follow by differentiating the displayed phases, then
substituting their stationary equations into the second derivatives.
For bounded \(\sigma\), with \(a_t=\sqrt{t/(2\pi)}\), Taylor expansion
of the square root gives
\[
 w_+=a_t-\frac{\ii\sigma}{4\pi a_t}+O(t^{-3/2}),\qquad
 w_-=a_t+\frac{\ii(1-\sigma)}{4\pi a_t}+O(t^{-3/2}).
\]
On the line they are conjugates. Changing \(\sigma\) moves them
smoothly, and the chosen phase Hessians remain nonzero throughout the
high critical strip. The saddle definition is tied to this phase/amplitude
split; it is not an invariant zero detector for an arbitrary integrand
factorization.

If a hypothetical off-line zero \(s_0\) exists, the exact conclusion is
\(\mathscr R(s_0)=-\chi(s_0)\overline{\mathscr R(1-\bar s_0)}\).
After contour displacement this is cancellation of the finite sums,
gamma factors and contour remainders together. It imposes no extra
stationary equation on \(w_\pm\). Saddles can be individually
nondegenerate while their complete contributions cancel. No zero-free
claim can be deduced by inspecting only the leading saddle or by
discarding a remainder near a small value of the total expression.

\subsection{Exact evolution of a complete Riemann--Siegel formula}
Riemann--Siegel is not restricted to \(\lambda=0\). That is the slice
where the classical zeta formula and the prime baseline apply directly.
For \(\lambda>0\), the complete source proves the exact identity
\begin{equation}\label{eq:imaginary-Gaussian}
 H_\lambda(z)=\frac1{\sqrt\pi}\int_\R
       \ee^{-v^2}H_0(z-2\ii\sqrt\lambda\,v)\dd v.
\end{equation}
For \(T>0\), forward heat similarly gives
\begin{equation}\label{eq:real-Gaussian}
 H_{-T}(z)=\frac1{\sqrt\pi}\int_\R
       \ee^{-v^2}H_0(z+2\sqrt T\,v)\dd v.
\end{equation}
To prove both, substitute \eqref{eq:H}, interchange integrals, and use
\(\pi^{-1/2}\int\ee^{-v^2+2av}\dd v=\ee^{a^2}\).
For the imaginary shifts Tonelli's absolute majorant is
\(\int\Phi(u)\ee^{\lambda u^2+|\Im z||u|}\dd u<\infty\).
For real shifts, use the same bound with \(\lambda=0\); the oscillatory
Gaussian integral supplies \(\ee^{-Tu^2}\). Thus these are actual
identities, not formal application of an unbounded reverse heat
operator to arbitrary data.

Write \(M(s)=s(s-1)\pi^{-s/2}\Gamma(s/2)/16\), so that
\(H_0(z)=M(s(z))\zeta(s(z))\). Under the shift in
\eqref{eq:imaginary-Gaussian}, \(s(z)\) becomes \(s(z)-\sqrt\lambda v\).
Consequently a complete finite Riemann--Siegel decomposition can be
evolved by Gaussian integration of its arithmetic terms and its
remainder, provided their individual integrability or a justified
combined continuation is retained. Polymath explicitly carries this
out in the heat-evolved decomposition
\cite[Section 4, equations (35)--(41)]{polymath2019heat}.

The arithmetic weight is visible even before treating the gamma factor:
\begin{equation}\label{eq:heat-Dirichlet-weight}
 \frac1{\sqrt\pi}\int_\R\ee^{-v^2}
       n^{-s+\sqrt\lambda v}\dd v
 =n^{-s}\ee^{\lambda(\log n)^2/4}.
\end{equation}
This is why heat-evolved Riemann--Siegel formulas contain
\(b_n^\lambda=\ee^{\lambda(\log n)^2/4}\). For the actual completed
field the moving \(M(s-\sqrt\lambda v)\), reflected contribution,
cutoff and errors remain. The displayed scalar identity does not
justify multiplying the old cosine formula termwise by \(b_n^\lambda\)
and leaving everything else unchanged. In particular the corresponding
infinite weighted Dirichlet series is not asserted to converge for
positive \(\lambda\).

The published explicit approximation is stated in a definite region:
\(0<\lambda\le1/2\), \(0\le\Im z\le1\), \(\Re z\ge200\), with
gamma corrections, shifted arguments, finite cutoff and error bounds
\cite[Theorem 1.3]{polymath2019heat}. Its exact decomposition has a wider
scope than that asymptotic theorem. A Gaussian average samples points
outside a bounded strip, so a strip-only error estimate cannot be
integrated over the whole Gaussian line without a separate tail payment.
Neither these identities nor the heat equation assign the time-zero
Euler product to \(H_\lambda\) for \(\lambda\ne0\).

\subsection{From heat-density asymptotics to a logarithmic error}
The heat-density investigation uses the same complete source with
forward time \(T=-\lambda\). Its positive-time estimates and its
time-zero Riemann--Siegel estimates have different regimes. For example,
a remainder proved on \(1/\log X\le T\le2\) cannot be specialized to
\(T=0\); the separate time-zero formula must be used. The positive
Newman-time formulas and crossover bounds likewise retain their stated
parameter bands. The source-use and transfer note records these domain
distinctions \cite{GNPHDATransfer}.

At zero Newman time a useful quantitative bridge is available. For
real \(x>0\), put \(s_\pm=1/2\pm\ii x/2\),
\(a_x=\sqrt{x/(4\pi)}\), \(N_x=\lfloor a_x\rfloor\), and
\(\vartheta_x=1-2(a_x-N_x)\). Let
\[
 U_x=\exp\{-\ii[(x/4)\log(x/(4\pi))-x/4-\pi/8]\}.
\]
With the coefficient convention \(C_k\) of the cited expansion, take
\begin{equation}\label{eq:RS-approximant}
 f_K(x)=2\Re\left\{M(s_+)\left[
 \sum_{n\le N_x}n^{-s_+}+(-1)^{N_x-1}U_xa_x^{-1/2}
             \sum_{k=0}^KC_k(\vartheta_x,1/2)a_x^{-k}\right]\right\}.
\end{equation}
The imported explicit remainder, for every \(x>0\) and integer \(K\ge1\), is
\begin{equation}\label{eq:RS-remainder}
 |H_0(x)-f_K(x)|\le\frac{2^{7/4}}7|M(s_+)|
       \Gamma((K+1)/2)1.1^{K+1}a_x^{-K-3/2}.
\end{equation}
This is the normalization recorded and checked in the inherited
transfer, using Polymath's Proposition 6.2
\cite{GNPHDATransfer,polymath2019heat}; the coefficients are those of
that exact version, not a new recurrence asserted here.

The gamma scale cancels in the logarithmic consumer. Indeed
\(D(x)M(s_-)=q(s_-)\), \(|q(s_-)|=1\) and
\(|M(s_+)|=|M(s_-)|\). Thus on the unit interval \(I_v\), with
\(a_- =\sqrt{(v-1/2)/(4\pi)}\), \(D f_K\) approximates \(F\) with
the explicit uniform error
\begin{equation}\label{eq:RS-normalized-error}
 \delta_{v,K}=\frac{2^{7/4}}7
       \Gamma((K+1)/2)1.1^{K+1}a_-^{-K-3/2}.
\end{equation}
The approximant need not be zero-free or holomorphic in a whole disk
for Corollary~\ref{cor:canonical-log}; holomorphy is required of the
exact field. Changes of the integer cutoff do not authorize omitting
the compensating remainder.

The distinction between fixed and increasing expansion order is
quantitative. In \eqref{eq:local-canonical-constants},
\(n_v\sim\log v/\kappa\) and the prefactor of the envelope is
\(O(\log^2v)\). For \(K=2p\), use
\(\Gamma(p+1/2)\le\sqrt\pi\,p^p\) to obtain
\[
 \log(2\delta_{v,2p}/M_v)
 \le-(p+7/4)\log v+O(p\log p).
\]
Choosing \(p=\lceil(4/\kappa)\log\log v\rceil\) therefore gives
\begin{equation}\label{eq:RS-log-transfer}
 \mathcal E_v(\delta_{v,2p})
 =O(\log^2v\,\ee^{-\kappa p+o(1)})=O((\log v)^{-2})
 \qquad(v\to\infty).
\end{equation}
This proves that an increasing-order, complete Riemann--Siegel error
can pass through the deteriorating local logarithmic exponent. A
fixed algebraic source-error rate does not make this particular
sufficient envelope vanish. At fixed height the remainder hierarchy
is asymptotic: the ratio of the bounds at orders \(2p+2\) and \(2p\)
is \((p+1/2)1.1^2/a_-^2\). Sending \(p\to\infty\) at fixed \(v\)
is therefore not a convergence claim. The finite initial-height range
requires its own acquisition, and the arithmetic sum still contains
order \(\sqrt v\) modes.

Equation~\eqref{eq:RS-log-transfer} controls an error allowance. The
remaining time-zero arithmetic task is to bound the signed logarithm of the
retained expression against \(P(X)-G(X,1)\) plus the chosen defect margin, including exceptional
centers. A saddle description, or a remainder tending to zero, does
not supply that comparison. This is the precise boundary between the
background established here and the transition-scale development below.
The time-zero gamma normalization and its pole correction must be
distinguished from the harmonic carrier used there.
